\captionsetup[longtable]{justification=raggedright,singlelinecheck=false}

\begin{longtable}{ | p{0.1\textwidth} | p{0.30\textwidth} | p{0.50\textwidth} |}
\caption{Kebutuhan Fungsional Sistem}
\label{tbl:fungsional} \\
\hline
\multicolumn{1}{|p{0.1\textwidth}|}{\centering\textbf{Kode}} & \multicolumn{1}{p{0.30\textwidth}|}{\centering\textbf{Parameter}} & \multicolumn{1}{p{0.50\textwidth}|}{\centering\textbf{Deskripsi Spesifikasi}} \\
\hline
\endfirsthead

\multicolumn{3}{l}{Tabel \ref{tbl:fungsional} Kebutuhan Fungsional Sistem (lanjutan)} \\
\hline
\multicolumn{1}{|p{0.1\textwidth}|}{\centering\textbf{Kode}} & \multicolumn{1}{p{0.30\textwidth}|}{\centering\textbf{Parameter}} & \multicolumn{1}{p{0.50\textwidth}|}{\centering\textbf{Deskripsi Spesifikasi}} \\
\hline
\endhead


\hline
\multicolumn{3}{r}{\textit{bersambung ke halaman berikutnya}} \\
\endfoot

\hline
\endlastfoot

F-01 & \textit{Structured Object Representation} & Sistem harus mengubah skema objek ke dalam suatu representasi terstruktur yang membedakan setiap komponen, \textit{field}, dan relasi antar-objek sehingga dapat diakses kembali secara terprogram pada tahap \textit{retrieval}. \\
\hline

F-02 & \textit{Structured Relation Retrieval} & Sistem harus menyediakan mekanisme \textit{retrieval} yang mampu mengambil himpunan objek dan \textit{field} yang saling berkaitan (misalnya hubungan antara \textit{Deployment}, \textit{Pod}, dan \textit{Service}) berdasarkan entitas awal dan \textit{intent} pengguna sebagaimana didefinisikan pada F-03. \\
\hline

F-03 & \textit{Intent Classification} & Sistem harus mampu (1) mendeteksi entitas Kubernetes yang dirujuk pengguna (misalnya \textit{Service}, \textit{Pod}), dan (2) mengklasifikasikan maksud pengguna ke dalam beberapa kategori kueri yang berbeda sehingga strategi \textit{retrieval} dan format keluaran dapat disesuaikan dengan karakteristik tiap jenis kueri. \\
\hline

F-04 & \textit{Context Construction \& Injection} & Sistem harus mengubah hasil F-02 menjadi konteks terstruktur yang dapat digunakan sebagai basis generasi jawaban. \\
\hline

F-05 & \textit{Semantic Embedding} & Sistem harus menghasilkan representasi vektor (\textit{embedding}) untuk setiap node dalam representasi terstruktur dan menyimpannya ke dalam indeks vektor, sehingga pencarian berbasis kemiripan semantik dapat digunakan sebagai mekanisme \textit{fallback} pada tahap \textit{retrieval}. \\
\hline

F-06 & \textit{Web Interaction Interface} &
Sistem harus menyediakan antarmuka web interaktif yang memungkinkan pengguna memasukkan pertanyaan dalam bahasa alami, menampilkan jawaban beserta konfigurasi YAML (jika relevan), serta menyajikan jejak penalaran (\textit{reasoning path}) secara visual disertai pesan kesalahan yang informatif. \\
\hline

F-07 & \textit{Conditional YAML Generation Constraint} & Sistem harus secara kondisional menghasilkan keluaran berupa blok kode YAML yang valid secara sintaksis apabila pengguna meminta pembuatan konfigurasi, dan menghasilkan jawaban naratif untuk jenis kueri lainnya yang tidak bersifat generatif. \\
\hline


\end{longtable}
